% GENERATED FILE. Edit canonical lessons, then run npm run generate:books.
% canonical-source-hash: fnv1a64:3f4d576a49a85bed
% canonical-lessons: JA-C118-nawa, JA-C118-hashigo, JA-C118-kusari, JA-C118-wata, JA-C118-honya

\chapter{Rope, Ladder, Chain, Cotton wool, Bookshop}
\label{ch:ja-rope-ladder-chain-cotton-wool-bookshop}
\hlchaptermodality{118}

\begin{chapteropening}
\textbf{By the end of this chapter:} \emph{I can say a rope, a ladder, a chain, cotton wool and a bookshop.}

Complete the last lesson of chapter 118: i can say a rope, a ladder, a chain, cotton wool and a bookshop.
\end{chapteropening}

\section[\texorpdfstring{\ja{なわ} (nawa)}{なわ (nawa)}]{\ja{なわ} (nawa) — a rope (of straw)}

\label{lesson:JA-C118-nawa}



\begin{quote}
Before the new one: say the Japanese for a Japanese serow, then the Japanese for a water flask.
\end{quote}

\subsection*{You'll want to know: \ja{なわ}}
\textbf{\ja{なわ}} — \emph{nawa} — \textquotedblleft{}a rope (of straw)\textquotedblright{}.

\begin{grammarlens}[title={What you've built}]
One more word for this chapter.
\end{grammarlens}

\subsection*{Your turn}
\begin{itemize}
\raggedright
  \item \emph{Say it:} \emph{nawa}
  \item \emph{Say it:} \emph{nawa}, once more
  \item \emph{Say it:} say \emph{nawa}
  \item \emph{Recall:} say the Japanese for a Japanese serow, then the Japanese for a water flask, then say \emph{nawa} again
\end{itemize}

\begin{tcolorbox}[breakable,colback=teal!4,colframe=teal!35!black,title={Before you move on}]
What does \ja{なわ} mean? (\textquotedblleft{}A rope (of straw)\textquotedblright{}.) Say it once more.
\end{tcolorbox}
\section[\texorpdfstring{\ja{はしご} (hashigo)}{はしご (hashigo)}]{\ja{はしご} (hashigo) — a ladder}

\label{lesson:JA-C118-hashigo}



\begin{quote}
Before the new one: say the Japanese for a water flask, then the Japanese for a rope.
\end{quote}

\subsection*{You'll want to know: \ja{はしご}}
\textbf{\ja{はしご}} — \emph{hashigo} — \textquotedblleft{}a ladder\textquotedblright{}.

\begin{grammarlens}[title={What you've built}]
One more word for this chapter.
\end{grammarlens}

\subsection*{Your turn}
\begin{itemize}
\raggedright
  \item \emph{Say it:} \emph{hashigo}
  \item \emph{Say it:} \emph{hashigo}, once more
  \item \emph{Say it:} say \emph{hashigo}
  \item \emph{Recall:} say the Japanese for a water flask, then the Japanese for a rope, then say \emph{hashigo} again
\end{itemize}

\begin{tcolorbox}[breakable,colback=teal!4,colframe=teal!35!black,title={Before you move on}]
What does \ja{はしご} mean? (\textquotedblleft{}A ladder\textquotedblright{}.) Say it once more.
\end{tcolorbox}
\section[\texorpdfstring{\ja{くさり} (kusari)}{くさり (kusari)}]{\ja{くさり} (kusari) — a chain}

\label{lesson:JA-C118-kusari}



\begin{quote}
Before the new one: say the Japanese for a rope, then the Japanese for a ladder.
\end{quote}

\subsection*{You'll want to know: \ja{くさり}}
\textbf{\ja{くさり}} — \emph{kusari} — \textquotedblleft{}a chain\textquotedblright{}.

\begin{grammarlens}[title={What you've built}]
One more word for this chapter.
\end{grammarlens}

\subsection*{Your turn}
\begin{itemize}
\raggedright
  \item \emph{Say it:} \emph{kusari}
  \item \emph{Say it:} \emph{kusari}, once more
  \item \emph{Say it:} say \emph{kusari}
  \item \emph{Recall:} say the Japanese for a rope, then the Japanese for a ladder, then say \emph{kusari} again
\end{itemize}

\begin{tcolorbox}[breakable,colback=teal!4,colframe=teal!35!black,title={Before you move on}]
What does \ja{くさり} mean? (\textquotedblleft{}A chain\textquotedblright{}.) Say it once more.
\end{tcolorbox}
\section[\texorpdfstring{\ja{わた} (wata)}{わた (wata)}]{\ja{わた} (wata) — cotton wool}

\label{lesson:JA-C118-wata}



\begin{quote}
Before the new one: say the Japanese for a ladder, then the Japanese for a chain.
\end{quote}

\subsection*{You'll want to know: \ja{わた}}
\textbf{\ja{わた}} — \emph{wata} — \textquotedblleft{}cotton wool\textquotedblright{}.

\begin{grammarlens}[title={What you've built}]
One more word for this chapter.
\end{grammarlens}

\subsection*{Your turn}
\begin{itemize}
\raggedright
  \item \emph{Say it:} \emph{wata}
  \item \emph{Say it:} \emph{wata}, once more
  \item \emph{Say it:} say \emph{wata}
  \item \emph{Recall:} say the Japanese for a ladder, then the Japanese for a chain, then say \emph{wata} again
\end{itemize}

\begin{tcolorbox}[breakable,colback=teal!4,colframe=teal!35!black,title={Before you move on}]
What does \ja{わた} mean? (\textquotedblleft{}Cotton wool\textquotedblright{}.) Say it once more.
\end{tcolorbox}
\section[\texorpdfstring{\ja{ほんや} (honya)}{ほんや (honya)}]{\ja{ほんや} (honya) — a bookshop}

\label{lesson:JA-C118-honya}



\begin{quote}
Before the new one: say the Japanese for a chain, then the Japanese for cotton wool.
\end{quote}

\subsection*{You'll want to know: \ja{ほんや}}
\textbf{\ja{ほんや}} — \emph{honya} — \textquotedblleft{}a bookshop\textquotedblright{}.

\begin{grammarlens}[title={What you've built}]
One more word for this chapter.
\end{grammarlens}

\subsection*{Your turn}
\begin{itemize}
\raggedright
  \item \emph{Say it:} \emph{honya}
  \item \emph{Say it:} \emph{honya}, once more
  \item \emph{Say it:} say \emph{honya}
  \item \emph{Recall:} say the Japanese for a chain, then the Japanese for cotton wool, then say \emph{honya} again
\end{itemize}

\begin{tcolorbox}[breakable,colback=teal!4,colframe=teal!35!black,title={Before you move on}]
What does \ja{ほんや} mean? (\textquotedblleft{}A bookshop\textquotedblright{}.) Say it once more.
\end{tcolorbox}
